\documentclass[11pt]{article}

\usepackage[margin=1in]{geometry}
\usepackage{tabularx}
\usepackage{booktabs}
\usepackage{enumitem}
\usepackage{float}

\setlength{\emergencystretch}{3em}

\input{../Comments.text}
\input{../Common.text}

\title{Problem Statement and Goals\\\progname}
\author{\authname}
\date{}

\begin{document}

\maketitle

Table~\ref{TblRevisionHistory} records the changes to this document.

\begin{table}[H]
\caption{Revision History} \label{TblRevisionHistory}
\begin{tabularx}{\textwidth}{llX}
\toprule
\textbf{Date} & \textbf{Developer(s)} & \textbf{Change}\\
\midrule
Sept 22, 2026 & Justin Kwinecki & Initial problem statement and goals\\
Sept 26, 2026 & Justin Kwinecki & Finished Section 1\\
Sept 27, 2026 & Zayed Umer & Revise problem statement and goals;
see disclosure appendix\\
Sept 27, 2026 & Zayed Umer & Cut repetition and clarify wording\\
Sept 27, 2026 & Zayed Umer & Revise the problem section from review\\
Oct 11, 2026 & Justin Kwinecki & Pivot idea update \\
Oct 11, 2026 & Justin Kwinecki & Align scope and goals with semantic routing \\
Oct 14, 2026 & Justin Kwinecki & Clarify serving-engine integration focus \\
\bottomrule
\end{tabularx}
\end{table}

\section{Problem Statement}

\progname{} aims to reduce response latency when semantic routing switches
between models by reducing repeated prompt prefill for shared conversation
context.

\subsection{Problem}

Semantic routing selects a model for each request based on the request's
meaning or needs. In a multi-turn interaction, a router may send successive
requests with the same conversation history to different models. Before a
model can generate a response, it processes the prompt in a phase called
prefill. A newly selected model may need to prefill the conversation history
again, adding delay before it can respond.

Repeated prefill increases time to first token and uses computing
capacity that could serve other requests. Existing serving software can
reuse a prefix after it has been processed for the same model
\cite{vLLMPrefixCaching}; after a model switch, the selected model may need
to process that shared prefix again. This makes model-switch latency and
repeated prompt processing the project's measurable serving concerns. The
problem is to reduce that repeated work while preserving the selected
model's response and the independence of each request.

\subsection{Scope}
\label{sec:scope}

Tandem targets text-generation requests in a multi-turn interaction where
semantic routing may select different models for successive turns. The
requests share conversation history, while the newest user message may
differ. The project investigates serving-engine integration to reduce
repeated prefill on a model switch. Its planned approach is to transfer a
cross-model key-value (KV) cache into the serving engine's KV memory and
use it when compatible with the request. When reuse is unavailable or
unusable, the engine uses its native prefill path. The evaluation will use
one server and a defined set of models and serving configurations. It will
include same-model turns and turns routed to a different model, with and
without reusable context. If feasibility work requires a different reuse
approach, the same workload and native-prefill comparison will apply.

Evaluation will compare identical workloads, routing decisions, models,
cache hits, and generation settings with and without Tandem's changes. The
baseline will retain the serving engine's existing prefix reuse features.
The project excludes training the base language models, development of a
semantic router, agent orchestration, document retrieval, application
interfaces, and operation across multiple servers. Reuse of similar but
differently worded context is outside the initial scope. Reuse across
different models will be limited to model combinations for which correctness
can be established.

The project focuses on serving-engine integration and evaluation, rather
than developing a new method to predict one model's cache from another.
Evaluation will compare time to first token after a model switch using the
chosen reuse approach against the engine's native full-prefill path, while
checking output correctness and recording reuse failures and fallback
behavior. The
\href{../DevelopmentPlan/DevelopmentPlan.pdf}{Development Plan} describes
the implementation and evaluation approach.

\subsection{Inputs and Outputs}

The inputs are:
\begin{itemize}[noitemsep]
  \item Requests, each with an identifier, conversation history, latest
  user message, and generation settings such as maximum response length
  and sampling controls.
  \item The set of trained models available to the router and the encoding
  rules that convert text into tokens, the units each model processes.
  \item A routing decision for each request, specifying the selected model.
  \item For the planned transfer approach, a candidate KV cache and its
  source model, target model, associated conversation context, and
  compatibility metadata.
  \item For evaluation, request arrival times, the serving configuration,
  and a baseline run using the same routing decisions.
\end{itemize}

The outputs are:
\begin{itemize}[noitemsep]
  \item Generated text for each accepted request, with its identifier,
  selected model, and completion status or failure indication.
  \item Evaluation results: time to first token, end-to-end response time,
  completed requests per second, prompt tokens processed on each model,
  and comparison with the baseline for the same selected model. The
  comparison reports generated-text agreement and output-score differences
  under controlled settings. Results include the workload and configuration
  needed to reproduce them.
\end{itemize}

\subsection{Stakeholders}
\label{sec:stakeholders}

\begin{itemize}
  \item \textbf{Application developers, including independent developers:}
  need faster responses without changing their agents' tasks or results.
  \item \textbf{Serving-engine maintainers:} need changes they can
  integrate and maintain, supported by correctness and compatibility tests.
  \item \textbf{Infrastructure operators and deploying organizations:}
  need to serve more requests within their hardware capacity and budget.
  \item \textbf{Application end users:} need timely responses that are
  unaffected by other agents' private instructions.
  \item \textbf{The capstone team, TA, and instructor:} need a feasible
  scope and reproducible evidence that the project meets its goals.
\end{itemize}

\subsection{Environment}

Developers and operators will run Tandem on a local or remote Linux server
with a supported graphics processing unit (GPU) and enough memory for the
models and requests. They supply the models, routing decisions, and
compatible serving software. Applications submit requests through a
programmatic interface; end users interact with those applications.

Initial evaluation uses one server and a fixed set of model configurations
at a time. The Development Plan covers developer tools and laboratory
hardware.

\section{Goals}

These goals will be refined into acceptance criteria in the Software
Requirements Specification and test procedures in the Verification and
Validation Plan. Performance comparisons use the baseline in
Section~\ref{sec:scope}, repeated runs, and reported variation. Improvements
must exceed that variation.

\begin{enumerate}[label=G\arabic*., ref=G\arabic*]
  \item \textbf{Reduce repeated prefill after model switches.} Reduce the
  processing of shared conversation history when a request is routed to a
  different model. Compare the chosen reuse approach with the serving
  engine's native full-prefill path, and report time to first token,
  repeated prompt tokens, and reuse or fallback outcomes.
  \item \textbf{Increase serving capacity without unacceptable delay.}
  Seek higher completed requests per second on the same hardware and matched
  workloads, while reporting median and 95th-percentile time to first token
  alongside throughput.
  \item \textbf{Preserve model behaviour and request independence.}
  Compare generated text and output scores with the selected model's native
  serving baseline under controlled settings, including multi-token
  generation. Assess whether one request's private instructions affect
  another request's output scores.
  \item \textbf{Maintain ordinary serving.} Preserve the serving engine's
  ability to handle same-model requests and requests without shared context,
  and measure any effect of the integration on their delay and throughput.
\end{enumerate}

Evaluation will report improvements, regressions, and workloads with no
measurable change.

\section{Stretch Goals}

After meeting the core goals, the team may pursue:
\begin{itemize}
  \item \textbf{Support deeper branching.} Extend correctness and
  performance evaluation to conversations with several shared branching
  points, including branches introduced during an ongoing interaction.
  \item \textbf{Support more configurations.} Repeat the evaluation on
  additional models or GPUs to test whether the benefits persist.
  \item \textbf{Reduce memory demand.} Use less peak GPU memory than the
  matched baseline for shared-context workloads and measure whether this
  permits more concurrent requests within the same memory limit.
\end{itemize}

\section{Custom Documents (CDs)}

The Development Plan draft proposes these custom documents, subject to
instructor approval:

\begin{itemize}
  \item \textbf{Fall: Machine Learning Model Report.} Justify the selected
  model, evaluation prompts, and correctness checks. Discuss model and data
  limitations, and explain why training procedures do not apply.
  \item \textbf{Winter: Performance Report.} Analyze repeated processing,
  response delay, throughput, and resource use. Record configurations,
  workloads, repeated trials, variation, and failure cases. The Verification
  and Validation Report will reference this analysis when assessing
  whether requirements were met.
\end{itemize}

\bibliographystyle{plain}
\bibliography{../../refs/References}

\section*{Appendix --- Generative AI Use Disclosure}

OpenAI Codex revised all sections using Justin Kwinecki's initial draft
and proposed additions, the Development Plan draft, repository checklists,
and course outline. It checked the cited prefix-caching documentation,
edited for clarity and concision, and built and inspected the PDF. No
performance experiments were run or novelty established.

Zayed Umer supplied reflection notes about collaborative brainstorming,
the difficulty of finding a novel idea, and the motivation to improve
memory use even modestly. Codex edited those notes and connected them to
the proposed scope. The team remains responsible for reviewing the wording,
claims, and scope before submission.

Justin Kwinecki used ChatGPT (free web version) to research which groups
of professionals would benefit from this project and their specific use
cases, for Section~\ref{sec:stakeholders}, Stakeholders. The team reviewed
the resulting list before including it.

\section*{Appendix --- Reflection}

\begin{enumerate}
  \item \textbf{What went well while writing this deliverable?}

  We agreed early that we wanted to work on inference. That shared
  interest made brainstorming collaborative: we could build on each
  other's ideas and focus on which problem to tackle.

  \item \textbf{What pain points did you experience during this
  deliverable, and how did you resolve them?}

  Finding a novel idea was harder than choosing the topic. With so much
  inference research underway, we struggled to tell whether an idea
  addressed an open problem. That uncertainty remains. We plan to review
  existing work and measure the selected serving engine while testing the
  feasibility of cross-model cache reuse. The evaluation can adjust the
  reuse approach if feasibility work shows it is needed.

  \item \textbf{How did you and your team adjust the scope of your goals
  to ensure they are suitable for a Capstone project?}

  We first wanted to make inference use fewer resources, including memory.
  To keep the project manageable, we narrowed that interest to repeated
  prefill after model switches. We will judge success by lower time to
  first token and more completed requests per second, while preserving the
  selected model's behaviour and keeping requests independent. Same-model
  requests and requests without shared context must still work normally.

  Memory savings and support for more model configurations remain stretch
  goals. We can meet the core goals by improving response time and serving
  capacity without reducing peak GPU memory use. Any claim of memory
  savings will need a separate comparison with the matched baseline.
\end{enumerate}

\end{document}
